% ============================================================
% CAPÍTULO 2: DATOS
% ============================================================

\chapter{Datos}
\label{chap:datos}

La adquisición de datos de calidad es fundamental para el desarrollo de modelos capaces
de predecir con generalidad y fiabilidad. Por ello se han incorporado varias bases de datos
obtenidas de fuentes públicas, en concreto, de los repositorios de PhysioNet.
Estas bases de datos contienen registros de electrocardiogramas (ECG) con anotaciones
sobre episodios de fibrilación auricular paroxística (PAF).

\section{Bases de datos utilizadas}
\label{sec:bases_datos}

Para el desarrollo de este proyecto se han empleado 5 bases de datos detalladas a continuación:

\subsection{PAF Prediction Challenge Database (afpdb)}
\label{subsec:afpdb}
Esta base de datos~\cite{moody_paf_2001} fue diseñada para el PhysioNet Challenge de predicción de
Fibrilación Auricular Paroxística. Contiene registros de ECG de dos canales con
una duración de 30 minutos y una frecuencia de muestreo de 128~Hz. El conjunto de
aprendizaje consta de 50 parejas de registros (100 registros en total) procedentes
de 48 sujetos:
\begin{itemize}
  \item \textbf{Registros tipo 'n':} 50 registros procedentes de sujetos sanos o sin antecedentes de Fibrilación Auricular.
  \item \textbf{Registros tipo 'p':} 50 registros obtenidos inmediatamente antes de un episodio de Fibrilación Auricular (Pre-PAF).
\end{itemize}
Adicionalmente, existe un conjunto de test con otros 100 registros procedentes de 50 sujetos, que utilizaremos para compararnos
con otros trabajos publicados.

\subsection{China Physiological Signal Challenge 2021 (cpsc2021)}
\label{subsec:cpsc2021}
Esta base de datos~\cite{PhysioNet-cpsc2021-1.0.0} contiene 1425 registros de 105 sujetos,
dos canales, muestreadas a 200~Hz. Contiene anotaciones
de ritmo, en concreto, de inicio y fin de episodios de Fibrilación Auricular.

\subsection{Long-Term Atrial Fibrillation Database (ltafdb)}
\label{subsec:ltafdb}
Esta base de datos~\cite{petrutiu_long-term_2003} contiene registros de ECG ambulatorio de
larga duración (aproximadamente 24 horas por registro), grabados a una frecuencia de muestreo
de 128~Hz. Consta de 84 registros de 84 sujetos con fibrilación auricular
paroxística o sostenida.

\subsection{Smart Health Devices - Atrial Fibrillation Database (shdb-af)}
\label{subsec:shdb}
Esta base de datos~\cite{PhysioNet-shdb-af-1.0.1} proporciona señales de ECG de dos
canales registradas mediante dispositivos wearables en condiciones cotidianas.
Consta de 128 registros
de larga duración (24 horas) procedentes de 122 sujetos con PAF,
muestreados originalmente a 125~Hz y resampleados a 200~Hz.
Nos proporciona la asignación entre señal y sujeto para prevenir filtraciones (leakage)

\subsection{MIT-BIH Atrial Fibrillation Database (afdb)}
\label{subsec:afdb}
Esta base de datos~\cite{moody_mit-bih_1992} consta de 25 registros de ECG de dos
canales con una duración de 10 horas cada uno, grabados a 250~Hz.
Cabe destacar que de 25, 23 de estos registros contienen señales de
ECG utilizables, ya que los otros 2 constan únicamente de anotaciones.

\section{Características comparativas de los conjuntos de datos}
\label{sec:caracteristicas_datos}

A continuación se presenta una tabla~\ref{tab:resumen_datos} que resume las características de las 5 bases de datos utilizadas
en este proyecto, incluyendo el número de pacientes, el número de registros, el número
de canales y la frecuencia de muestreo de cada base de datos.

\begin{table}[htbp]
  \centering
  \caption{Resumen comparativo de las 5 bases de datos de estudio}
  \label{tab:resumen_datos}
  \begin{tabular}{lcccc}
    \toprule
    \textbf{Base de datos} & \textbf{Nº Pacientes} & \textbf{Nº Registros} & \textbf{Canales} & \textbf{Frec. Muestreo} \\
    \midrule
    \texttt{afpdb~\cite{moody_paf_2001}} & 48 (98) & 100 (300) & 2 & 128 Hz \\
    \texttt{cpsc2021~\cite{PhysioNet-cpsc2021-1.0.0}} & 105 & 1425 & 2 & 200 Hz \\
    \texttt{ltafdb~\cite{petrutiu_long-term_2003}} & 84 & 84 & 2 & 128 Hz \\
    \texttt{shdb-af~\cite{PhysioNet-shdb-af-1.0.1}} & 122 & 128 & 2 & 200 Hz \\
    \texttt{afdb~\cite{moody_mit-bih_1992}} & 25 & 25 (23) & 2 & 250 Hz \\
    \bottomrule
  \end{tabular}
  \vskip 5pt
  \small
  \textit{Nota:} Para \texttt{afpdb}, entre paréntesis se indica el total incluyendo el conjunto de test. Para \texttt{afdb}, entre paréntesis se indica el número de registros que contienen señales utilizables.
\end{table}
